\section{Building intuition: accumulation and random mixing}
\label{sec:spin-intuition}

Two simple instances explain the roles of the outer, interleaver, and inner.
Both use independent uniform linear injections
$C_i:\F_2^{B/2}\to\F_2^B$, where $B$ is even and $N=MB$ for $M$ blocks.
The outer concatenates their outputs, and a uniform permutation interleaves
the resulting $N$ bits. Setup is sampled once and fixed for every message.
For each fixed nonzero input to an outer block, its image is uniform over
the nonzero $B$-bit vectors. After interleaving, a word of weight $w$ has
uniform support among the $N$ positions. Thus the remaining question is
how the inner transforms such a support.

\paragraph{The minimal inner: an accumulator.}
For interleaved input $u$, set $y_0:=0$ and compute
\[
  y_i:=y_{i-1}+u_i,\qquad i=1,\ldots,N.
\]
Write $\Acc_N(u):=y$. This is the one-stage Block-Accumulate
architecture~\cite{KolesnikovEtAl2026BA}.
The inner has one bit of state. Each input one flips the output, so output
ones occupy the intervals between alternating flips. Counting these
intervals gives an exact input--output weight enumerator.

Let $U_w$ be uniform among $N$-bit vectors of weight $w$.
For $0<\delta<1/2$, the enumerator implies
\[
 \Pr[\wt(\Acc_N(U_w))\le\lfloor\delta N\rfloor]
 \le \bigl(2\sqrt{\delta(1-\delta)}\bigr)^w.
\]
The outer spectrum counts how many messages can supply each weight;
this bound controls their chance of producing a low-weight output.
Multiplying and summing gives the first moment from
Section~\ref{sec:framework}. Logarithmically growing outer blocks suffice
to make that sum vanish at a positive relative distance.

Theorem~\ref{thm:independent-accumulator-spin} makes this concrete:
with its block schedule $B\ge3\log_2N$, the rate-$1/2$ family has relative
distance greater than $0.0037$ with probability $1-o(1)$ over setup.
This modest guarantee already shows that a one-bit inner can suffice.
Together with a terminated streaming realization of the outer, it gives
a positive answer to the logarithmic-outer-memory question of
Bazzi, Mahdian, and Spielman~\cite{bms09}; Appendix~\ref{sec:bms-question}
states the precise interpretation. The minimal construction is useful
independently of the stronger mixing considered next.

\paragraph{A stronger inner: random recursive convolution.}
Keep the same outer and interleaver, but use the recursive inner from
Expand--Convolute~\cite{RaghuramanRindalTanguy2023EC}. For memory $m$, sample
independent $\alpha_i\gets\F_2^m$ at setup, independently of the outer and
interleaver, and set $y_j:=0$ for $j\le0$. Compute
\[
 y_i:=u_i+\langle\alpha_i,(y_{i-1},\ldots,y_{i-m})\rangle.
\]
For a fixed input, condition on the preceding outputs. If the state is
nonzero, the fresh coefficient vector makes the next output bit uniform.
If the state is zero, the next output equals the input bit. Tracking the
number of trailing output zeros, clipped at $m$, therefore gives an exact
finite-state calculation of the expected weight enumerator. The zero-state
case matters: the outputs are not an independent uniform string.

Theorem~\ref{thm:independent-random-spin} combines this calculation with
logarithmic outer block size and memory. At rate $1/2$, it proves relative
distance greater than $0.11002$ with probability $1-o(1)$, close to the
Gilbert--Varshamov value of approximately $0.110028$. The full transfer and
the bounds for sparse and linear weights appear together in
Appendix~\ref{sec:random-spin}.

\paragraph{Toward efficient encoding.}
Both direct encoders require $\Theta(N\log N)$ work under their proved
schedules: the outer blocks are dense, and the random convolution evaluates
logarithmically many taps per coordinate. The next section replaces these
costly components with spectrum-controlled constituents, structured routing,
and batched inner operations. The proof must then account for the remaining
coordinate structure and the inner's actual state law. Random SPIN motivates
that design; its distance theorem does not by itself establish the
structured construction's guarantee.
